\StartChapter{Komplexe Zahlen\ZSFScriptRef{14,18,19}}

\SubsectionBar{Eigenschaften}

Eine \ZSFkeyword{komplexe Zahl} ist von der Form $z = a + ib$ wobei $a = \Re(z)$ der \ZSFkeyword{Realteil} und $b = \Im(z)$ der \ZSFkeyword{Imaginärteil} ist.

Die \ZSFkeyword{komplex konjugierte Zahl} $\bar{z} = a - ib$ ist die Spiegelung von $z$ an der reellen Achse (meist die $x$-Achse).

\textVorBox{Es gelten die folgenden Regeln für Konjugation und Betrag:}

\begin{tablebox}
\begin{ZSFtable}{C C}
    \ZSFheaderRow \ZSFhead{Konjugation} & \ZSFhead{Betrag} \\
    $\displaystyle \overline{z_1 + z_2} = \overline{z_1} + \overline{z_2}$ & $\displaystyle |z_1 z_2| = |z_1||z_2|$ \\
    $\displaystyle \overline{z_1 z_2} = \overline{z_1} \, \overline{z_2}$ & $\displaystyle \left| \frac{z_1}{z_2} \right| = \frac{|z_1|}{|z_2|}$ \\
    $\displaystyle \overline{z^n} = \bar{z}^n$ & $\displaystyle |z_1 + z_2| \leq |z_1| + |z_2|$ \\
    \ZSFspan{C}{2}{$\displaystyle z\bar{z} = |z|^2$} \\
\end{ZSFtable}
\end{tablebox}
\ZSFindex{Betrag}
\ZSFindex{Dreiecksungleichung}
\ZSFindexsee{Konjugierte}{komplex konjugierte Zahl}

Die Darstellung $z = r(\cos(\varphi) + i\sin(\varphi))$ nennt man \ZSFkeyword{Polarform}.
\ZSFindexsee{Polardarstellung}{Polarform}

\textVorBox{Die Darstellung $z = re^{i\varphi}$ nennt man \ZSFkeyword{Eulersche Form}. Es gilt:}
\ZSFindexsee{Euler-Formel}{Eulersche Form}
\begin{formulabox}
    \ZSFformulaline{r = |z| = \sqrt{a^2 + b^2}}{Betrag}
\end{formulabox}

\textVorBox{und:}
\begin{tablebox}
\begin{ZSFtable}{C C C}
    \ZSFheaderRow \ZSFhead{$a > 0$} & \ZSFhead{$a < 0,\ b \geq 0$} & \ZSFhead{$a < 0,\ b < 0$} \\
    $\displaystyle \varphi = \arctan\left(\frac{b}{a}\right)$ & $\displaystyle \varphi = \arctan\left(\frac{b}{a}\right) + \pi$ & $\displaystyle \varphi = \arctan\left(\frac{b}{a}\right) - \pi$ \\
\end{ZSFtable}
\end{tablebox}
\textNachBox{wobei $\varphi$ gegen den Uhrzeigersinn gemessen wird.}
\ZSFindex{Argument}
\ZSFindexsee{Phasenwinkel}{Argument}%
\textVorBox{\ZSFdanger{Wichtig:}}
\begin{formulabox}
    $\displaystyle z = e^{i\pi} = -1,\quad z = e^{i\frac{\pi}{2}} = i$
\end{formulabox}

\SubsectionBar{Rechenregeln}

\textVorBox{Die Eulersche Form eignet sich für komplizierte Operationen mit komplexen Zahlen:}

\begin{tablebox}
\begin{ZSFtable}[header=false]{Q{0.7} Z{1.3}}
     Multiplikation: & $\displaystyle z_1 \cdot z_2 = \ZSFmhlA{r_1 \cdot r_2} \cdot \ZSFmhlB{e^{i(\varphi_1 + \varphi_2)}}$ \\
    Division: & $\displaystyle \frac{z_1}{z_2} = \ZSFmhlA{\frac{r_1}{r_2}} \ZSFmhlB{e^{i(\varphi_1 - \varphi_2)}}$ \\
     Inverse: & $\displaystyle \frac{1}{z} = \ZSFmhlA{\frac{1}{r}} \ZSFmhlB{e^{-i\varphi}}$ \\
    Potenz: & $\displaystyle z^n = \ZSFmhlA{r^n} \ZSFmhlB{e^{in\varphi}}$ \\
     Wurzel: & $\displaystyle \sqrt[n]{z} = \ZSFmhlA{\sqrt[n]{r}} \ZSFmhlB{e^{i\left(\frac{\varphi + 2k\pi}{n}\right)}}$ \\
\end{ZSFtable}
\end{tablebox}
\textNachBox{wobei $k = 0, 1, 2, \dots, n-1$}
\textNachBox{Alle Wurzeln liegen auf Kreis mit Radius $\sqrt[n]{r}$}
\ZSFindex{Wurzel (komplex)}

\ZSFNewColumn
\SubsectionBar{Polynome\ZSFScriptRef{21,57}}

\textVorBox{Ein \ZSFkeyword{Polynom} ist eine Summe von Potenzen einer Variable:}
\begin{formulabox}
    \ZSFformulaline{p(x) = ax^n + bx^{n-1} + \dots + cx + d}{Polynom}
\end{formulabox}
\textNachBox{$x$ wird als \ZSFkeyword{Nullstelle} des Polynoms $p$ bezeichnet, wenn $p(x) = 0$.}

\begin{ZSFlist}
    \ZSFItem{Grad}{Grösster Exponent eines Polynoms (oben $n$).}
    \ZSFItem{Anzahl Nullstellen}{Ein Polynom vom Grad $n$ hat $n$ Nullstellen.}
    \ZSFItem{Vielfachheit}{Eine Nullstelle kann einfach ($p=0$), doppelt ($p=0$, $p'=0$), dreifach ($p=p'=p''=0$), bis $n$-fach vorkommen.}
    \ZSFItem{Komplexe Nullstellen}{Treten paarweise auf ($z,\bar{z}$). Eine einfach komplexe Nullstelle entspricht zwei, eine doppelt komplexe vier Nullstellen. Ein Polynom mit ungeradem Grad besitzt mindestens eine reelle Nullstelle.}
\end{ZSFlist}
\ZSFindex{Vielfachheit}
\ZSFindex{Grad eines Polynoms}

\textVorBox{Es gelten folgende Regeln für den Grad eines Polynoms:}
\begin{formulabox}
    \[\begin{aligned}
        \Grad\{p + q\} &\leq \max(\Grad\{p\}, \Grad\{q\}) \\
        \Grad\{pq\} &= \Grad\{p\} + \Grad\{q\}
    \end{aligned}\]
\end{formulabox}
